\begin{examplebox}
\textbf{例1（右支：直接求角平分线交点与切线）。}

双曲线 $\dfrac{x^2}{9}-\dfrac{y^2}{16}=1$，
焦点 $F_1(-5,0),F_2(5,0)$，取点
$P\left(\dfrac{15}{4},3\right)$。
设 $\angle F_1PF_2$ 的内角平分线交 $x$ 轴于 $M(m,0)$。
求 $m$ 及 $PM$ 方程。

\textbf{解：} $a=3,b=4,c=5,x_0=15/4$，直接套用
\[
m=\frac{a^2}{x_0}=\frac{9}{15/4}=\boxed{\frac{12}{5}},
\qquad\boxed{M\left(\frac{12}{5},0\right)}.
\]
对应切线为
\[
\frac{(15/4)x}{9}-\frac{3y}{16}=1
\quad\Longleftrightarrow\quad
\boxed{20x-9y=48}.
\]
\textbf{校验：} $PF_1=37/4,\ PF_2=13/4$，
\[
\frac{F_1M}{MF_2}
=\frac{5+12/5}{5-12/5}=\frac{37}{13}
=\frac{PF_1}{PF_2}.
\]

\medskip
\textbf{例2（左支：完全不用另背一套公式）。}

在同一双曲线上取 $Q\left(-\dfrac{15}{4},3\right)$，
其焦点三角形的顶角内角平分线交轴于 $N(n,0)$。
直接得
\[
\boxed{n=\frac{9}{-15/4}=-\frac{12}{5}},
\qquad\boxed{-20x-9y=48}.
\]
左支的 $n$ 为负，且 $-a<n<0$，与几何位置完全一致。

\medskip
\textbf{例3（反向设点：先知轴截距，再求曲线上点）。}

设 $a=3,b=4$，双曲线上一点 $P(x_0,y_0)$ 的焦点三角形
顶角内角平分线交 $x$ 轴于 $M\left(\dfrac32,0\right)$。
求所有可能的 $P$。

\textbf{解：}从最简结论反推
\[
x_0=\frac{a^2}{m}=\frac{9}{3/2}=6.
\]
代入双曲线
\[
\frac{36}{9}-\frac{y_0^2}{16}=1
\quad\Longrightarrow\quad y_0^2=48.
\]
所以
\[
\boxed{P=(6,4\sqrt3)\quad\text{或}\quad(6,-4\sqrt3)}.
\]
两点都不是顶点，故均满足题设。

\medskip
\textbf{例4（内切圆切点：分清 $T$ 和 $M$）。}

仍取例1的 $P\left(\frac{15}{4},3\right)$，
作 $\triangle PF_1F_2$ 内切圆，设其与 $x$ 轴切于 $T$。
因为 $P$ 在右支，
\[
\boxed{T=(a,0)=(3,0)},\qquad
\boxed{M=\left(\frac{12}{5},0\right)}.
\]
故 $T\ne M$。进一步，内切圆半径
\[
\rho=\frac{[PF_1F_2]}{s}
=\frac{\frac12\cdot10\cdot3}{\frac12(37/4+13/4+10)}
=\boxed{\frac43},
\]
内心为 $I=(3,4/3)$（点 $P$ 在 $x$ 轴上方）。
这也检验了 $P,I,M$ 三点共线。

\medskip
\textbf{例5（考场变式：切线横截距的范围）。}

设 $P$ 在双曲线的右支且不是顶点，
$PM$ 为顶角内角平分线。因为 $x_0>a$，
\[
0<m=\frac{a^2}{x_0}<a.
\]
所以交点 $M$ 必落在\textbf{原点与右顶点之间}，
并且随着 $P$ 的横坐标变大，$m$ 单调变小。
左支则有 $-a<m<0$。
\end{examplebox}
